\section{Internet Control Message Protocol (ICMP)}

Protocolo helper de IP completamente \textbf{informativo}, es decir, no toma decisiones sino que permite reportar errores o consultar información. Tampoco agrega confiabilidad, sino que \textbf{brinda feedback}. Se encapsula en IP, pero no es un protocolo de transporte.


Cuando un dispositivo descarta un datagrama IP, genera un mensaje ICMP que incluye el encabezado (header) del datagrama original que produjo el problema y al menos los primeros 8 bytes del campo de datos de dicho datagrama.

Los mensajes ICMP son los siguientes.

\subsection{Mensajes de Error (Error-reporting messages)}
Estos mensajes reportan errores que un dispositivo puede encontrar cuando procesa un paquete IP.

\begin{itemize}[label=\textcolor{acento}{$\bullet$}]
    \item \textbf{Destination Unreachable (Tipo 3):} Indica que el paquete no pudo llegar a su destino.
    \begin{itemize}[label=\textcolor{acento!60}{$\bullet$}]
        \item \textit{Code 0:} Net Unreachable (Red inalcanzable).
        \item \textit{Code 1:} Host Unreachable (Host inalcanzable).
        \item \textit{Code 3:} Port Unreachable (Puerto inalcanzable, muy común al fallar una conexión UDP).
    \end{itemize}
    
    \item \textbf{Time Exceeded (Tipo 11):} Se genera cuando expira un temporizador relacionado al datagrama.
    \begin{itemize}[label=\textcolor{acento!60}{$\bullet$}]
        \item \textit{Code 0:} In Transit. Se venció el campo TTL (Time to Live) durante el tránsito por los routers.
        \item \textit{Code 1:} Reassembly Timeout. Se descartó un datagrama fragmentado porque se agotó el tiempo de espera antes de que llegaran todos sus fragmentos para el reensamblado.
    \end{itemize}
\end{itemize}

\subsection{Mensajes de Consulta (Query messages)}
Estos mensajes ocurren usualmente de a pares (una solicitud y su correspondiente respuesta) y permiten obtener información específica del estado de un dispositivo en la red.

\begin{itemize}[label=\textcolor{acento}{$\bullet$}]
    \item \textbf{Echo Request / Echo Reply (Tipos 8 y 0):} Son los mensajes utilizados por la clásica utilidad de diagnóstico \texttt{ping} para comprobar la conectividad bidireccional entre nodos.
    \item \textbf{Timestamp Request / Timestamp Reply (Tipos 13 y 14):} Utilizados para determinar el tiempo de tránsito a través de la red o para sincronizar relojes de los dispositivos.
\end{itemize}


Cuando un mensaje ICMP es descartado, \textbf{no se genera otro} para avisar que se descartó (evito bucles). Solo avisamos el descarte de un fragmento si es el \textbf{primero}, y tampoco avisamos si se descarta un mensaje con destino multicast.